\section{A2 EEG 语义融合模型}

这一章介绍我们课程设计中用到的主要定量模型——A2 temporal-spectral-spatial EEG 语义融合模型。我们做这个模型的出发点很简单：把 EEG 当成视觉语义增强信号来用。具体来说，当输入图像退化、质量变差时，我们用同步采集的 EEG 信号来补充和修正退化图像在 CLIP 语义空间里的表示，从而让语义识别更准确。这个模型的全称是 Temporal-Spectral-Spatial Transformer，我们简称为 A2。它的做法是在 CLIP 语义空间里同时提取 EEG 的三个维度特征——时间、频谱和空间通道，然后通过一个门控残差融合机制来补充退化图像的语义表征。

\subsection{设计动机}

EEG 本质上是一个多维时间序列信号。我们用的数据里，每条 EEG trial 的尺寸是 $[64, 250]$，也就是 64 个电极通道在 250 个时间点上的电压值。这里面包含三个不同维度的信息：时间维度记录了图像出现之后不同时刻的脑电活动，比如早期视觉诱发电位（P100/N170）和晚期的语义加工成分（P300）；频谱维度反映了各个频段（$\alpha$、$\beta$、$\gamma$）的能量振荡分布，很多研究都表明这些频段跟视觉感知、注意力和语义加工有密切关系；空间维度描述了 64 个电极在头皮上的位置分布，反映了不同脑区之间的功能连接关系。

如果只用单一类型的编码器，通常只能利用其中一到两个维度的信息。举个例子，纯时序卷积能很好地抓住局部时间模式，但完全不管电极之间的空间关系；频谱 CNN 从时频域提取特征，但丢失了原始的时间拓扑信息，而且很难建模跨频段的交互；Transformer 直接处理原始时间序列可以捕获长程依赖关系，但缺少显式的频谱建模。所以我们设计 A2 的思路是：对三个维度分别设计专门的编码分支，每个分支用适合该维度特点的结构，最后通过后融合层来做跨维度的特征交互。在我们模型容量有限的情况下，这种先分开提取再融合的做法，比单纯把单一编码器加深要高效得多。

在具体结构选择上，我们参考了 EEG 信号处理领域已有的研究成果。ConvTransformer 架构已经被验证在 raw EEG 时间序列建模上效果不错，因为卷积层能提供局部时间平滑，Transformer 层能捕获全局时序依赖。基于 STFT 的频谱 CNN 也已经在人脑 EEG 情感识别和运动想象分类任务上有了成功应用。通道空间注意力机制在 EEG 脑机接口里被用来学习电极之间的功能连接。A2 做的新事情就是把这三个方向整合到一个端到端可训练的框架里，专门针对视觉语义解码这个任务来做优化。

\subsection{A2 EEG Encoder 结构}

A2 encoder 的结构是三个并行分支加上一个后融合 Transformer。图~\ref{fig:a2-encoder} 画出了整体结构。

\placeholderfigure{A2 temporal-spectral-spatial EEG encoder 结构图}{a2-encoder-placeholder}{5.0cm}

\textbf{Temporal 分支}：我们用了多尺度 ConvTransformer 结构。具体做法是用四个不同卷积核大小的 Conv1d（$k \in \{3, 7, 15, 31\}$）并行提取多尺度时间特征，然后把各条支路的输出拼起来，经过 $1\times1$ 卷积做混合，再通过 4 层 Transformer Encoder 来建模长程时序依赖。最后用 learnable attention pooling 加 MLP 投影得到 512 维的 temporal embedding。多尺度设计的好处是既能捕捉毫秒级的尖峰变化，也能关注到跨越几百毫秒的慢变趋势。

\textbf{Spectral 分支}：这是一个基于可微分短时傅里叶变换的频谱 CNN。我们先对 EEG 信号做 STFT（窗长 64、hop length 16），对幅度谱取对数然后逐通道做标准化，再用两层 2D CNN 提取时频域的二维模式。最后经过全局平均池化和 MLP 投影，得到 512 维的 spectral embedding。

\textbf{Spatial 分支}：我们把 EEG 信号转置一下，让 64 个通道各自的时间序列都通过线性层投影成 hidden\_dim 维的 token，然后用 learnable attention pooling 做加权聚合，再经过 MLP 输出 512 维的 spatial embedding。这个分支的好处是不需要预先定义通道图结构，靠自注意力机制自己学会电极之间的功能连接模式。

\textbf{后融合}：三个分支各自输出 512 维向量，拼起来变成 1536 维，然后经过 2 层 Transformer Encoder 做跨分支的交互融合。这个融合 Transformer 用了 learnable branch type embedding 来标记每个特征来自哪个分支。最后通过 attention pooling 和 MLP 投影器得到 512 维的融合 EEG embedding $f_{\mathrm{eeg}}$。

\subsection{CLIP 文本原型与语义融合}

CLIP 文本原型沿用第 2 章的做法：对 40 个类别分别构造 prompt ``a photo of a \{class\_name\}''，用冻结的 CLIP text encoder 得到归一化的文本原型向量。在做语义预测的时候，我们把 A2 输出的融合 EEG embedding $f_{\mathrm{eeg}}$ 和退化图像的视觉 embedding $f_{\mathrm{img}}$ 通过 gated residual fusion 融合在一起：

\begin{equation}
  f_{\mathrm{fused}} = f_{\mathrm{img}} + \sigma\left(g\left([f_{\mathrm{img}}; f_{\mathrm{eeg}}]\right)\right) \odot \Delta(f_{\mathrm{eeg}})
\end{equation}

这里面 $g(\cdot)$ 是门控网络（一个 2 层 MLP），$\sigma$ 是 sigmoid 函数，$\Delta(\cdot)$ 是 EEG 残差变换网络。门控机制的作用是：根据图像和 EEG 的联合信息，动态控制 EEG 注入的强度。如果图像质量比较好，门控系数就接近零，EEG 的干预就少一些；如果图像退化比较严重，门控系数就变大，EEG 信息以更大权重注入进来。

融合之后得到的向量 $f_{\mathrm{fused}}$ 跟 40 个类别的文本原型算余弦相似度，再通过 softmax 得到每个类别的预测概率。图~\ref{fig:a2-fusion} 展示了这个流程。

\placeholderfigure{A2 semantic fusion 流程图}{a2-fusion-placeholder}{4.0cm}

\subsection{训练设计}

训练 A2 语义融合模型的时候，我们主要用交叉熵分类损失作为目标。给定融合嵌入 $f_{\mathrm{fused}}$ 和 40 个类别的文本原型 $\{p_c\}_{c=1}^{40}$，分类预测是通过对余弦相似度做 softmax 得到的：

\begin{equation}
  \mathcal{L}_{\mathrm{CE}} = -\frac{1}{B} \sum_{i=1}^{B} \log \frac{\exp(\tau^{-1} \cos(f_{\mathrm{fused},i}, p_{y_i}))}{\sum_{c=1}^{K} \exp(\tau^{-1} \cos(f_{\mathrm{fused},i}, p_c))}
\end{equation}

这里的 $\tau$ 是可学习的温度参数，$y_i$ 是第 $i$ 个样本的真实类别标签。

我们的模型更新参数只在真实配对的 EEG 数据上进行。Shuffled EEG 和 random EEG 只在测试时作为对照来评估，不参与训练。这个训练和评估分离的做法是我们实验设计的关键：如果 shuffled/random EEG 也拿来训练，模型可能学到的是 EEG 通道里的任意噪声模式，而不是真正的语义信息。只有 real EEG 参与训练，然后用全部四种 EEG 模式来评估，我们才能保证模型学到的 EEG 利用能力确实来自真实配对的神经信号。

训练用的是 AdamW 优化器，初始学习率设成 $1\times10^{-4}$，权重衰减是 0.01。批大小是 64（在单 GPU 上通过梯度累积来实现），用了 cosine 学习率衰减策略，最低学习率降到初始值的 $1/100$。为了降低小样本情况下的过拟合风险，我们把 EEG encoder 里的 dropout rate 设为 0.2，同时在 gated fusion MLP 里加了 LayerNorm 正则化。总共训练 80 个 epoch，在验证集上用 top-1 accuracy 来做早停选择，patience 设成 25 个 epoch。

A2 主语义融合模型的核心训练目标就是交叉熵分类损失。EEG-to-CLIP 的 InfoNCE/MSE 对齐训练主要是用在前期对齐实验或者可选的预训练阶段，并不作为最终 A2 语义融合模型的唯一训练目标。在我们典型的训练流程里，第一阶段是可选的 EEG-to-CLIP 预对齐：在冻结的 CLIP 语义空间中，用 InfoNCE 对比损失把 EEG embedding 拉向对应图像的 CLIP embedding。这样预对齐之后，后续的融合训练就有了比较好的初始化。第二阶段是语义融合训练（这是核心阶段）：加载对齐阶段的 checkpoint（或者从随机初始化开始），联合训练 gated fusion 模块，用交叉熵分类损失来优化。在第二阶段里，CLIP ViT 和 CLIP text encoder 都是完全冻结的，只更新 A2 EEG encoder、gated fusion 和温度参数 $\tau$。本文第 8 章报告的所有 A2 结果都是以交叉熵损失为主训练目标，配合可选的 InfoNCE 预对齐来训练的。

\subsection{模型验证方法}

为了客观判断 EEG 是不是真的包含了机器可以利用的语义信息，我们在评估 A2 的时候设置了多重对照：real EEG 是被试观看图像时的同步 EEG 信号；shuffled EEG 是把同被试内的 EEG trial 跟图像随机打乱配对，这样做会破坏时间同步关系，但保留了 EEG 本身的统计分布；random EEG 是从高斯分布采样出来的，完全不包含任何真实的神经语义信息。只有 real EEG 系统地超过了 shuffled 和 random EEG，我们才能说 EEG 提供了有效的语义辅助。这个对照思路贯穿了我们课程设计的所有实验。

\subsection{小结}

A2 temporal-spectral-spatial EEG encoder 的主要特点可以总结成这样：三个并行分支分别提取 EEG 的时间、频谱和空间信息；gated residual fusion 根据输入质量动态调节 EEG 的辅助强度；CLIP 文本原型提供了一个受约束的语义分类空间，避免了在小样本上从头训练的麻烦；严格的 shuffled/random EEG 对照给我们提供了客观的效果验证。具体的实验结果在第 8 章详细给出。
